%==============================================================================%
% Chapter 4 – Formal Verification Framework
%==============================================================================%

\chapter{Formal Verification Framework}
\label{ch:formal_framework}

\section{Overview of the Verification Architecture}

The verification of the \haetae{} NTT is structured as a
\emph{vertical refinement chain}: a sequence of proof steps, each relating
an implementation or specification at one abstraction level to a
specification at the next.  \Cref{fig:arch} illustrates the overall structure.

\begin{figure}[htbp]
\centering
\begin{tikzpicture}[
  box/.style = {rectangle, draw, rounded corners, minimum width=8cm,
                minimum height=1.2cm, align=center, font=\small},
  arrow/.style = {-Stealth, thick},
  label/.style = {font=\footnotesize\itshape, midway, right=3pt}
]
  \node[box, fill=red!10]   (jasmin) at (0, 6)
        {\textbf{Layer 0: Jasmin Source} \\ \texttt{poly.jinc, reduce.jinc, hpoly.jazz}};

  \node[box, fill=orange!10] (extract) at (0, 4.2)
        {\textbf{Layer 1: Extracted EasyCrypt Module} \\ \texttt{Hpoly\_extract.ec}};

  \node[box, fill=yellow!10] (impspec) at (0, 2.4)
        {\textbf{Layer 2: Imperative EasyCrypt Spec} \\ \texttt{NTT\_Fq.ec, Fq.ec}};

  \node[box, fill=green!10]  (algspec) at (0, 0.6)
        {\textbf{Layer 3: Algebraic/Spectral Spec} \\
         \texttt{NTTFullSpec.ec, NTTFullAlgebra.ec}};

  \draw[arrow] (jasmin) -- (extract)
    node[label] {Jasmin compiler extraction};
  \draw[arrow] (extract) -- (impspec)
    node[label] {pRHL equivalence (RefJasminNTT.ec)};
  \draw[arrow] (impspec) -- (algspec)
    node[label] {Loop invariant / induction (NTTFullAlgebra.ec)};

  \node[box, fill=blue!10, minimum width=4cm, minimum height=0.9cm]
    (endtoend) at (5.5, 3.3)
    {\textbf{End-to-End} \\ \texttt{NTTEndToEnd.ec}};

  \draw[arrow, dashed, bend right=20] (jasmin.east) to
    node[right=3pt, font=\footnotesize\itshape] {top-level theorem}
    (endtoend.north);
  \draw[arrow, dashed, bend left=20] (algspec.east) to (endtoend.south);
\end{tikzpicture}
\caption{Four-layer verification architecture for the \haetae{} NTT.
         Solid arrows are machine-checked proof steps; dashed arrows indicate
         the end-to-end composition.}
\label{fig:arch}
\end{figure}

\section{Layer 0: Jasmin Source}

The \jasmin{} source files implement the NTT algorithm in a style close to
assembly but with structured control flow.  The key files are:
\begin{itemize}
  \item \texttt{poly.jinc}: The main NTT and inverse-NTT kernels,
        implemented as \jasm{inline fn \_poly\_ntt} and
        \jasm{inline fn \_poly\_invntt}.

  \item \texttt{reduce.jinc}: The Montgomery reduction primitive
        \jasm{\_\_montgomery\_reduce} and the derived \jasm{\_\_fqmul}.

  \item \texttt{zetas.jinc}: The 256-element twiddle-factor table as a
        global constant.

  \item \texttt{hpoly.jazz}: The exported entry points
        \jasm{export fn poly\_ntt\_jazz} and
        \jasm{export fn poly\_invntt\_jazz}.
\end{itemize}

The \jasmin{} compiler verifies safety properties (no array-out-of-bounds, no
use-before-initialisation) and then simultaneously generates:
\begin{enumerate}
  \item x86-64 assembly (\texttt{hpoly.s}).
  \item An \easycrypt{} module \texttt{Hpoly\_extract.ec} that is
        semantically equivalent to the assembly by the Jasmin compiler
        correctness theorem~\cite{Almeida2017}.
\end{enumerate}

\section{Layer 1: Extracted EasyCrypt Module}

The extracted module \texttt{Hpoly\_extract.ec} (\texttt{extract/} folder)
contains \easycrypt{} procedure definitions that mirror the \jasmin{}
structure.  A fragment is shown in \Cref{lst:extract}.

\begin{lstlisting}[language=EasyCrypt,
  caption={Fragment of extracted procedure in Hpoly\_extract.ec},
  label={lst:extract}]
module M = {
  proc poly_ntt_jazz (rp : BArray1024.t) :
      BArray1024.t = {
    var zeta : W32.t;
    var t    : W32.t;
    var len  : int;
    var start: int;
    var j    : int;
    (* ... body mirrors Jasmin structure ... *)
    return rp;
  }
}.
\end{lstlisting}

This module is the ``ground truth'' for what the compiled assembly computes.
All subsequent proof obligations are stated with respect to this module.

\section{Layer 2: Imperative EasyCrypt Specification}

The file \texttt{NTT\_Fq.ec} contains hand-written \easycrypt{} procedures
that closely mirror the reference C implementation.  Unlike the extracted
module, these procedures work with abstract coefficient arrays over \(\F_q\)
(type \ec{coeff Array256.t}) rather than machine-word arrays.

The translation from machine words to field coefficients is captured by the
\emph{word-to-coefficient} relation \ec{poly\_repr}: a concrete
\ec{BArray1024.t} word array corresponds to a coefficient array when each
loaded 32-bit word maps to the same \(\F_q\) element.  The strengthened
\ec{poly\_repr\_bound} predicate also records the signed bit-width bound
required for word arithmetic.

The imperative specification exposes the full loop structure, making it
amenable to Hoare-style invariant reasoning.

\section{Layer 3: Abstract Spectral Specification}

The files \texttt{NTTFullSpec.ec} and \texttt{NTTFullAlgebra.ec} define the
NTT at the level of abstract algebra.  No loop structure is visible; the
NTT is defined as a linear evaluation map over $\F_q^{256}$.

\begin{definition}[Abstract forward NTT, \ec{NTTFullSpec.full\_ntt}]
  Let $p : \F_q^{256}$.  The forward NTT of $p$ is the map
  \[
    \NTT(p)[i] = \sum_{j=0}^{255} p[j] \cdot \zeta_0^{(2 \brev_8(i)+1) \cdot j}
    \pmod{q},\quad i = 0, \dots, 255.
  \]
\label{def:full_ntt}
\end{definition}

\begin{definition}[Abstract inverse NTT, \ec{NTTFullSpec.full\_invntt}]
  Let $p : \F_q^{256}$.  The inverse NTT of $p$ is the map
  \[
    \INTT(p)[i] = \frac{1}{256} \sum_{j=0}^{255} p[j] \cdot
    \zeta_0^{-((2\brev_8(j)+1) \cdot i)} \pmod{q},\quad i = 0, \dots, 255.
  \]
\label{def:full_intt}
\end{definition}

\section{End-to-End Composition}

The file \texttt{NTTEndToEnd.ec} imports all three layers and states the
main correctness theorems as Hoare judgements for the exported \jasmin{}
procedures:

\begin{lstlisting}[language=EasyCrypt,
  caption={End-to-end theorem statement in NTTEndToEnd.ec},
  label={lst:e2e}]
lemma haetae_jasmin_ntt_correct p :
  hoare [Hpoly_extract.M.poly_ntt_jazz :
    NTT_Fq.poly_repr_bound rp p 16 ==>
    NTT_Fq.poly_repr_bound res (NTTFullSpec.full_ntt p) 24].

lemma haetae_jasmin_invntt_correct p :
  hoare [Hpoly_extract.M.poly_invntt_jazz :
    NTT_Fq.poly_repr_bound rp p 16 ==>
    NTT_Fq.poly_repr_bound res
      (NTT_Fq.array256_mont (NTTFullSpec.full_invntt p)) 16].
\end{lstlisting}

These theorem statements capture the vertical refinement: starting from a
\jasmin{} export whose input array satisfies \ec{poly\_repr\_bound}, the
forward output represents the abstract spectral NTT with a 24-bit bound, and
the inverse output represents the abstract inverse NTT in Montgomery form
with a 16-bit bound.

\section{Proof Methodology}

\subsection{Proceeding Bottom-Up}

We build proofs bottom-up, starting with the most concrete layer and
moving toward the most abstract.  This order is forced by the dependency
structure: the pRHL refinement proof (Layer 1 $\to$ Layer 2) requires
that the imperative specification (Layer 2) is already proved correct with
respect to the abstract specification (Layer 2 $\to$ Layer 3).

\subsection{Hoare Logic for Loop Invariants}

The main technical tool for Layers 2--3 is \emph{loop invariant reasoning}.
Each NTT stage is characterised by an invariant asserting that a
partial set of coefficients has been transformed.  The invariant is a
conjunction of:
\begin{enumerate}
  \item A \emph{partial NTT} condition: indices below the current
        counter have been transformed through the current stage.
  \item A \emph{bounds condition}: transformed coefficients remain
        within the established bit-width bounds.
  \item A \emph{twiddle-factor synchronisation} condition: the twiddle
        counter matches the current loop indices.
\end{enumerate}

\subsection{pRHL for Procedure Equivalence}

The refinement from the extracted module to the imperative specification
uses pRHL judgements of the form
\[
\begin{aligned}
  &\Pr[\ec{Hpoly\_extract.M.\_poly\_ntt(rp)} \to \ec{rp'}] \\
  &\quad =
    \Pr[\ec{NTT\_Fq.NTT.ntt(a, zetas)} \to \ec{a'}],
    \quad \text{given } \ec{poly\_repr\_bound(rp,a,16)}.
\end{aligned}
\]
where \ec{poly\_repr\_bound} relates the concrete word array \ec{rp} to the
abstract field array \ec{a} and records the signed bit-width bound.  Since
both procedures are deterministic, this reduces to a functional equivalence
assertion.

\subsection{Field-Theory Automation}

The abstract-level proofs rely on \easycrypt{}'s theory-cloning mechanism to
instantiate the standard field axioms for $\F_q$ and to invoke automated
tactics for polynomial identities.  Long algebraic chains are dispatched by
the \ec{field} and \ec{ring} tactics, backed by CVC5 and Z3 for non-linear
arithmetic goals.

\section{File Dependency Graph}

\begin{figure}[htbp]
\centering
\begin{tikzpicture}[
  scale=0.95, transform shape,
  node distance=1.35cm and 2.25cm,
  box/.style={rectangle, draw, rounded corners, font=\ttfamily\scriptsize,
              minimum width=3.25cm, minimum height=0.8cm, align=center},
  arr/.style={-Stealth}
]
\node[box] (fq)      {Fq.ec};
\node[box, right=of fq] (gfq)     {GFq.ec};
\node[box, right=of gfq] (rq)      {Rq.ec};

\node[box, below=of fq]  (nttfq)   {NTT\_Fq.ec};
\node[box, below=of gfq] (mont)    {Montgomery.ec};
\node[box, below=of rq]  (fe)      {Fastexp.ec};

\node[box, below=of nttfq] (fullspec) {NTTFullSpec.ec};
\node[box, below=of mont]  (fullalg)  {NTTFullAlgebra.ec};

\node[box, below=of fullspec] (refjasmin) {RefJasminNTT.ec};
\node[box, below=of fullalg]  (extract)   {Hpoly\_extract.ec};

\node[box, below right=1.2cm and -1cm of refjasmin] (e2e) {NTTEndToEnd.ec};

\draw[arr] (fq) -- (nttfq);
\draw[arr] (gfq) -- (nttfq);
\draw[arr] (mont) -- (fq);
\draw[arr] (fe) -- (nttfq);
\draw[arr] (rq) -- (fullspec);
\draw[arr] (gfq) -- (fullspec);
\draw[arr] (fullspec) -- (fullalg);
\draw[arr] (nttfq) -- (fullalg);
\draw[arr] (fq) -- (refjasmin);
\draw[arr] (nttfq) -- (refjasmin);
\draw[arr] (extract) -- (refjasmin);
\draw[arr] (refjasmin) -- (e2e);
\draw[arr] (fullalg) -- (e2e);
\draw[arr] (fullspec) -- (e2e);
\end{tikzpicture}
\caption{Core EasyCrypt file dependency graph for the \haetae{} NTT
         verification.  Arrows point from dependency to dependent; low-level
         Jasmin array clones and standard-library theories are omitted here
         and recorded in \Cref{tab:ec_file_relationships}.}
\label{fig:dep_graph}
\end{figure}

\section{Relationships Between the EasyCrypt Files}
\label{sec:ec_file_relationships}

\Cref{fig:dep_graph} gives the main proof-chain dependencies.  The table below
records the role of each project \texttt{.ec} file and the reason it appears in
the checked development.  This distinction matters because not every import has
the same proof meaning: some files provide representation infrastructure, some
state executable procedures, and the last files compose already-checked
judgements into the final theorem.

\small
\begin{longtable}{
  >{\raggedright\arraybackslash}p{0.28\textwidth}
  >{\raggedright\arraybackslash}p{0.22\textwidth}
  >{\raggedright\arraybackslash}p{0.39\textwidth}}
\caption{Relationships between the EasyCrypt source files}
\label{tab:ec_file_relationships}\\
\toprule
\textbf{File} & \textbf{Layer / role} & \textbf{Relationship to the proof} \\
\midrule
\endfirsthead
\toprule
\textbf{File} & \textbf{Layer / role} & \textbf{Relationship to the proof} \\
\midrule
\endhead
\bottomrule
\endfoot

\texttt{Array256.ec} &
Array infrastructure &
Clones the Jasmin fixed-size array theory at size 256.  It is the common
container for polynomials, twiddle tables, and abstract NTT outputs. \\

\texttt{BArray1024.ec} and \texttt{WArray1024.ec} &
Concrete memory infrastructure &
Instantiate byte-array and word-array theories used by the extracted Jasmin
module.  They support table reads, word loads, and stores in
\texttt{Hpoly\_extract.ec} and \texttt{RefJasminNTT.ec}. \\

\texttt{GFq.ec} &
Field foundation &
Defines the \haetae{} modulus \(q=64513\) and clones the finite-field theory.
It supplies the coefficient type used by \texttt{Rq.ec}, \texttt{NTT\_Fq.ec},
\texttt{NTTFullSpec.ec}, and \texttt{NTTFullAlgebra.ec}. \\

\texttt{Rq.ec} &
Polynomial-ring vocabulary &
Defines the polynomial type, ring operations, root \(\zeta\), bit-reversal
notation, and big-sum infrastructure.  Its older \texttt{ntt} operator is kept
for comparison, while the final target is the full transform in
\texttt{NTTFullSpec.ec}. \\

\texttt{Montgomery.ec} &
Arithmetic foundation &
Contains modular-reduction lemmas and Montgomery arithmetic facts used to prove
the word-level routines in \texttt{Fq.ec} and reused by the Jasmin refinement
proofs. \\

\texttt{Fq.ec} &
Word arithmetic correctness &
Proves correctness of the extracted Montgomery reduction, field multiplication,
addition, subtraction, and signed-word bounds.  These lemmas discharge the
arithmetic side conditions inside \texttt{RefJasminNTT.ec}. \\

\texttt{Fastexp.ec} &
Exponentiation support &
Provides generic finite-field exponentiation lemmas.  \texttt{NTT\_Fq.ec}
clones this theory for the \haetae{} coefficient field. \\

\texttt{Hpoly\_extract.ec} &
Extracted Jasmin model &
Contains the EasyCrypt procedures generated from the Jasmin implementation:
\texttt{poly\_ntt\_jazz}, \texttt{poly\_invntt\_jazz}, their internal cores,
and the concrete twiddle tables.  It is the concrete side of the refinement
proof. \\

\texttt{NTT\_Fq.ec} &
Imperative specification &
Defines the hand-written coefficient-level procedures
\texttt{NTT.ntt} and \texttt{NTT.invntt}, the mathematical twiddle tables, and
the representation predicates such as \texttt{poly\_repr\_bound}.  It is the
target of \texttt{RefJasminNTT.ec} and the source of \texttt{NTTFullAlgebra.ec}. \\

\texttt{NTTFullSpec.ec} &
Closed-form abstract specification &
Defines \texttt{full\_ntt} and \texttt{full\_invntt} as the full 256-point
\haetae{} transforms.  It also records that \texttt{Rq.ntt} is not this full
transform, preventing the end-to-end theorem from targeting the wrong
operator. \\

\texttt{NTTFullAlgebra.ec} &
Imperative-to-abstract bridge &
Proves the Hoare statements that running \texttt{NTT\_Fq.NTT.ntt} with the
checked twiddle table returns \texttt{NTTFullSpec.full\_ntt}, and similarly for
the inverse transform.  This file closes the algebraic loop-invariant part of
the proof. \\

\texttt{RefJasminNTT.ec} &
Jasmin-to-imperative refinement &
Proves pRHL equivalences between the extracted Jasmin cores in
\texttt{Hpoly\_extract.ec} and the imperative procedures in \texttt{NTT\_Fq.ec}.
It also lifts the internal-core result to the public Jasmin wrappers. \\

\texttt{NTTEndToEnd.ec} &
Final composition &
Imports \texttt{RefJasminNTT.ec}, \texttt{NTTFullAlgebra.ec},
\texttt{NTTFullSpec.ec}, \texttt{NTT\_Fq.ec}, and
\texttt{Hpoly\_extract.ec}.  It composes the refinement and algebraic Hoare
lemmas into the final forward and inverse Jasmin NTT correctness statements. \\

\end{longtable}
\normalsize

In short, the proof has two independent bridges that meet in
\texttt{NTTEndToEnd.ec}.  The concrete bridge is
\[
  \texttt{Hpoly\_extract.ec}
  \xrightarrow{\ \texttt{RefJasminNTT.ec}\ }
  \texttt{NTT\_Fq.ec},
\]
which says that the extracted Jasmin procedures implement the imperative
coefficient-level model.  The algebraic bridge is
\[
  \texttt{NTT\_Fq.ec}
  \xrightarrow{\ \texttt{NTTFullAlgebra.ec}\ }
  \texttt{NTTFullSpec.ec},
\]
which says that the imperative model computes the full closed-form
\haetae{} NTT.  The end-to-end file imports both bridges and composes them.
